\documentclass[11pt]{article}
\usepackage[letterpaper,margin=0.85in]{geometry}
\usepackage[T1]{fontenc}
\usepackage{lmodern,amsmath,amssymb,amsthm,booktabs,longtable,array}
\usepackage{microtype,xcolor,enumitem,fancyhdr}
\usepackage[hidelinks]{hyperref}
\definecolor{ink}{HTML}{183D51}
\newtheorem{theorem}{Theorem}[section]
\newtheorem{lemma}[theorem]{Lemma}
\newtheorem{proposition}[theorem]{Proposition}
\newcommand{\Z}{\mathbb Z}
\newcommand{\N}{\mathbb N}
\newcommand{\Con}{\operatorname{Con}}
\newcommand{\ZFC}{\mathrm{ZFC}}
\newcommand{\SOS}{\operatorname{SOS}}
\newcommand{\code}[1]{\texttt{\small #1}}
\setlength{\parskip}{4pt}
\setlength{\parindent}{0pt}
\setlength{\emergencystretch}{2em}
\setlist{topsep=3pt,itemsep=2pt}
\pagestyle{fancy}\fancyhf{}
\fancyhead[L]{\small Symmetric residue budgets}
\fancyhead[R]{\small Diophantine consistency}
\fancyfoot[C]{\thepage}
\setlength{\headheight}{14pt}
\allowdisplaybreaks[3]
\hypersetup{pdftitle={Symmetric residue budgets for a 258-variable consistency quartic},pdfauthor={AI-assisted research prepared for Vladimir Reshetnikov and ProveIt}}
\begin{document}
\begin{center}
{\color{ink}\LARGE\bfseries Symmetric residue budgets\\[4pt]
for a 258-variable consistency quartic}\par
\vspace{8pt}
{\large A direct equation for $\Con(\ZFC)$, with compact companions}\par
\vspace{7pt}
Prepared for Vladimir Reshetnikov and the ProveIt repository\\
AI-assisted research contribution --- 25 October 2026
\end{center}

\begin{abstract}
We give a fixed polynomial $P_{258}\in\Z[x_0,\ldots,x_{257}]$ of exact
degree four such that ZFC proves that $P_{258}$ has no integer zero if and
only if ZFC is consistent. The equation is supplied completely as a sum
of 186 explicitly printed quadratic residual squares and as all 5,319
collected coefficients. Its coefficient height is 47,281. It refines the
latest inspected direct numerical-proof-calculus construction, which has
282 variables, 192 residuals, 5,491 terms and height 128,446.

The main change replaces eleven private interval certificates by shared
quadratic bounds on signed remainders. A radix estimate proved before the
factorial fixed point makes congruence uniqueness recover the missing
signs. Two adjacent scales can share a bound. A compatible odd-square
radix replaces a further sign certificate, while a changed recovery order
avoids circular reasoning. Thirty elementary residual-row operations
then reduce coefficient height without changing the zero set.

A companion quartic has 261 variables and 185 residuals. Two separate
ten-positive-variable equations save one arithmetic operation each, giving
77 operations at degree 630 and 78 at degree 468. Their proof removes an
auxiliary positivity shift by excluding a low Pell index. The new ordinary
proofs compose with explicitly retained repository calibrations. No global
simplicity minimum or new end-to-end proof-assistant certificate is claimed.
\end{abstract}

\section{The exact result and the comparison boundary}

The comparison snapshot is ProveIt~\cite{repo} commit
\code{e0c98c629bf10d30b695884f6b1d2ab361bdf513}.
We inspected the two requested research areas and all six ZIP archives
then present in \code{docs/incoming}. In particular, the comparison below
includes the incoming round-17 cofinal-anchor construction~\cite{cofinal}; it is not a
comparison against an obsolete earlier collection overview.

\begin{theorem}[The supplied direct quartic]
Let $r_0,\ldots,r_{185}\in\Z[x_0,\ldots,x_{257}]$ be the literal
quadratic polynomials in Appendix~\ref{app:residuals}, and define
\begin{equation}
 \boxed{\quad P_{258}(x_0,\ldots,x_{257})
     =\sum_{j=0}^{185}r_j(x_0,\ldots,x_{257})^2=0.\quad}
 \label{eq:main}
\end{equation}
All 258 coordinates range independently over $\Z$. Composing the proofs
in this report with the retained numerical C222/refutation calibration,
\begin{equation}
 \ZFC\vdash
 \left(\neg\exists\mathbf x\in\Z^{258}\;
       P_{258}(\mathbf x)=0\right)
       \longleftrightarrow\Con(\ZFC).
 \label{eq:consistency}
\end{equation}
The polynomial has exact degree four, 5,319 nonzero collected monomials,
and coefficient height 47,281. Every coordinate is live. Its coefficient
of $x_0^4$ is 13, and that of $x_{10}^2x_{12}^2$ is 47,281.
\end{theorem}

Here consistency means the absence of a finite syntactic refutation in
the repository's specified ZFC proof system. Equation~\eqref{eq:consistency}
is the result; it is not an independent assertion that either side is true.
The actual mathematical dependency is an ordinary proof formalizable in
ZFC. This packet does not supply a newly checked internal ZFC proof object
or a new complete Lean/Rocq build of the whole chain.

\begin{center}\small
\begin{tabular}{lrrrrr}
\toprule
Direct equation over $\Z$ & Variables & Rows & Degree & Terms & Height\\
\midrule
Incoming matched parent & 282 & 192 & 4 & 5,491 & 128,446\\
Shared four-budget companion & 261 & 185 & 4 & 5,332 & 47,281\\
\textbf{Main odd-square equation} & \textbf{258} & \textbf{186} & \textbf{4}
 & \textbf{5,319} & \textbf{47,281}\\
\bottomrule
\end{tabular}
\end{center}

The 258-variable result removes 24 variables and six rows from the matched
parent, removes 172 collected terms, and reduces its height by about
63.2\%. The 261-variable companion uses one fewer residual than the main
equation. Both have literal coefficient data and reconstruction checks.
These are achieved improvements in a specified direct lineage. Existing
native-encoding equations have fewer variables, sometimes at much higher
degree or with very large fixed coefficients. We do not call this a global
variable record or a globally simplest polynomial.

The direct branch encodes the numerical C222 proof calculus. It does not
simulate K215, specialize Jones's universal polynomial, or hide a machine
description in a large coefficient port. Its final coefficients are the
ordinary integers printed by its residuals. The separate compact branch
in Section~\ref{sec:compact-shift} retains its native/K215 ancestry and
must be assessed on its own domain and size metrics.

\section{The inherited arithmetic interface}

The parent direct construction is explicit and substantial. We retain its
complete guarded quadratic source, its source metadata, the ordinary
proofs of its arithmetic front, and the named logical interfaces. The
new proof changes a narrowly identified part of that source rather than
claiming to reprove the whole compilation from a finite test.

The logical chain is
\[
 \begin{aligned}
 \text{integer zero}&\ \Longleftrightarrow\ \text{finite C222 derivation}\\
 &\ \Longleftrightarrow\ \text{finite ZF refutation}\\
 &\ \Longleftrightarrow\ \text{finite ZFC refutation}.
 \end{aligned}
\]
The last two retained interfaces are:
\begin{itemize}
\item \code{KEndpoint.c222\_iff\_zf\_refutation};
\item \code{ConstructibleChoice.zf\_refutation\_iff\_zfc\_refutation}.
\end{itemize}
Their files are included as source evidence; inclusion is not a claim
that their entire transitive proof-assistant dependency tree was rebuilt.

The arithmetic construction packs finite local proof data by Chinese
remaindering, interprets binomial divisibilities as range conditions, and
uses explicitly bounded Pell certificates to represent powers. Its
signed orientation stage recovers both coordinates of each principal
Pell pair before it uses the low-digit remainder signs. This order is
the opportunity exploited here.

Write $t,F,B,x,y,V_0,q$ for its semantic arithmetic values and
$a_1,\ldots,a_9$ for its actual CRT representatives. They are polynomials
in the source graph's coordinates, not new variables in
\eqref{eq:main}. At the stage after power recovery but before the
factorial fixed point, the parent proves
\begin{gather}
 t\ge16,\quad F\ge1,\quad B\ge F+1\ge2,\qquad
 V_0=(F^2+1)^t,\qquad y=V_0(V_0-1),\label{eq:raw1}\\
 x\ge2+(t+1)y,\qquad a_i\ge0,\qquad
 H_0=1+\sum_{i=1}^9a_i\ge1,\qquad q=(BH_0)^x.
 \label{eq:raw2}
\end{gather}
In particular, this stage does not yet use $F=t!$. The raw ceiling is
$FB=V_0+s$ with $s\ge0$, and its two endpoint congruences give the large
lower bound on $x$. A modified order for the odd-square version will be
proved in Section~\ref{sec:odd}.

There are eleven modular numerator requests. Each has a denominator
$D_i=q^{k_i}$ and a low remainder $r_i$. The old rows impose
\begin{equation}
 D_i=r_i+s_i+1,\qquad
 r_is_i=u_i^2+v_i^2+w_i(w_i+1).
 \label{eq:oldguard}
\end{equation}
Once $D_i>0$ is known, the second row makes $r_i,s_i$ have the same sign,
and the first forces them both nonnegative. It costs three private integer
coordinates per request. We remove precisely the eleven product rows
in~\eqref{eq:oldguard}; all the first rows are retained. The three initial
ceiling/output sign graphs are separate from these eleven rows.

\section{A residue-recognition lemma with a shared bound}

The useful fact is not an unrestricted replacement for arbitrary interval
guards. It relies on the intended residues occupying a small part of their
interval.

\begin{lemma}[Symmetric residue recognition]\label{lem:recognition}
Let $D\ge1$ and let $\rho_1,\ldots,\rho_m$ be integers with
$0\le\rho_i<D/8$. Suppose integers $r_i$ satisfy
$r_i\equiv\rho_i\pmod D$ and
\begin{equation}
 D^2-4\sum_{i=1}^m r_i^2
      =1+u^2+v^2+w(w+1).
 \label{eq:budget}
\end{equation}
Then $r_i=\rho_i$ for all $i$. Conversely, if $m\le7$, the canonical
choice $r_i=\rho_i$ admits integers $u,v,w$ satisfying
\eqref{eq:budget}.
\end{lemma}
\begin{proof}
For every integer $w$, $w(w+1)\ge0$. Equation~\eqref{eq:budget}
implies $|r_i|<D/2$. Thus
\[
 |r_i-\rho_i|<D/2+D/8=5D/8<D.
\]
The difference is a multiple of $D$, so it is zero.

For the converse,
$4\sum_i\rho_i^2<7D^2/16<D^2$. Therefore the integer
$D^2-4\sum_i\rho_i^2-1$ is nonnegative. Every nonnegative integer $n$
has a representation $n=u^2+v^2+w(w+1)$. Indeed, the three-square theorem
applied to $4n+1$ gives exactly one odd coordinate and two even
coordinates; division by four gives the representation. This standard
ternary form is also used in the parent and in current formalized
Diophantine constructions~\cite{bayer}.
\end{proof}

The lemma allows the $r_i$ to be signed at the moment their congruences
are checked. Nonnegativity is a conclusion. Its applicability here rests
on proving the bound $\rho_i<D/8$ at the right point in the argument.

\section{The radix margin is available before factorial recovery}
\label{sec:margin}

The parent constructor ledger supplies the positive integer
$\ell=2^{22}=4,194,304$ and the derived-label bound
\[
 h\le(n_*+1)^\ell,
 \qquad n_*=\max(1,a,b,c,d),
\]
where $a,b,c,d$ are four of the nine nonnegative CRT representatives.
Since $n_*+1\le2H_0$ and $B\ge2$, we can use
\begin{equation}
 h\le(2H_0)^\ell\le(BH_0)^\ell.
 \label{eq:derivedbound}
\end{equation}
This deliberately permits $H_0=1$. We do not use the later proof that a
particular encoded index is nonzero.

By~\eqref{eq:raw1}, $V_0\ge2^t$, so
\[
 y\ge2^t(2^t-1)>\ell,\qquad y\ge t+2\qquad(t\ge16).
\]
Together with~\eqref{eq:raw2}, these imply the two coarse but sufficient
integer bounds
\begin{equation}
 x\ge\ell t+4,\qquad x\ge t(t+1)+4.
 \label{eq:gaps}
\end{equation}

\begin{lemma}[Uniform low-digit bound]\label{lem:digits}
For every one of the eleven numerator requests and every binomial
coefficient below its requested digit position,
\begin{equation}
 \binom nj\le q/16.
 \label{eq:smallcoeff}
\end{equation}
This conclusion follows from the pre-factorial facts
\eqref{eq:raw1}--\eqref{eq:raw2} and the constructor ledger.
\end{lemma}
\begin{proof}
There are three cases. For an ordinary representative $0\le n<H_0$,
the denominator exponent $k$ is at most $x$. If $0\le j<k$, then
\[
 \binom nj\le\max(1,n)^j\le H_0^x\le q/16,
\]
since $q=B^xH_0^x$, $B\ge2$ and $x\ge4$. This includes $n=0,j=0$.

For the derived-label representative, $k=t+1$ and $j\le t$.
Equations~\eqref{eq:derivedbound} and~\eqref{eq:gaps} give
\[
 \binom hj\le h^j\le(BH_0)^{\ell t}
       \le(BH_0)^{x-4}\le q/16.
\]

For the factorial upper index $n=F^2+t$, we have
$F^2+t\le B^{t+1}$: indeed $F^2<B^2\le B^t$,
$t\le2^t\le B^t$, and $2B^t\le B^{t+1}$. Therefore for $j<t$,
\[
 \binom{F^2+t}{j}\le(F^2+t)^t
      \le B^{t(t+1)}\le B^{x-4}\le q/16.
\]
None of the three arguments assumes $F=t!$.
\end{proof}

For $D=q^k$, define the canonical low residue in the proof by
\[
 \rho(n,k)=\sum_{j=0}^{k-1}\binom nj q^j.
\]
This is proof notation only; it is not a variable-length operation in the
final polynomial. The binomial theorem gives
$\rho(n,k)\equiv(q+1)^n\pmod D$, and Lemma~\ref{lem:digits} gives
\begin{equation}
 0\le\rho(n,k)\le\frac{q(D-1)}{16(q-1)}
       \le\frac{D-1}{8}<\frac D8.
 \label{eq:smallrho}
\end{equation}

The retained numerator equations make each actual $r_i$ congruent to
this canonical residue modulo its positive denominator. The factorial
numerator has a high part $D(B+qv)$; a range numerator has high part
$DNv$. Both are multiples of $D$ regardless of the still-unproved sign
of $v$. This is why the congruence can be used before the digit signs.

\section{Pooling, merging, and soundness of the direct graph}

The eleven requests have four denominator levels:
\begin{center}
\begin{tabular}{lrl}
\toprule
Denominator & Requests & Semantic upper indices\\
\midrule
$q^t$ & 2 & factorial upper index and encoded index\\
$q^{t+1}$ & 5 & four local inputs and derived label\\
$q^{y+1}$ & 2 & two local state representatives\\
$q^x$ & 2 & two strict beta quotient representatives\\
\bottomrule
\end{tabular}
\end{center}
For each level, introduce one row~\eqref{eq:budget} instead of its old
private product guards. The four rows use twelve guard coordinates in
place of thirty-three. After the inherited exact-power recovery,
Lemma~\ref{lem:recognition} and~\eqref{eq:smallrho} prove every
$r_i=\rho_i\ge0$. The retained linear rows then give
$s_i=D_i-1-r_i\ge0$. Thus every deleted old product guard can be restored
using a ternary representation of $r_is_i$.

All later inherited steps now apply. They authenticate the factorial
digit, obtain $F=t!$ from the ceiling fixed point, decode the amplified
CRT range conditions, and extract a finite C222 derivation from the
unit-weight proof-flow certificate. This proves soundness. For
completeness, start with the parent's witnesses for any finite refutation;
their canonical low residues satisfy~\eqref{eq:smallrho}, and the
converse of Lemma~\ref{lem:recognition} supplies the four new budgets.
The exact compiled result has 261 variables and 185 rows.

\subsection{Merging adjacent powers of the radix}

Let $D=q^t$ and $D'=qD=q^{t+1}$. Introduce just two product coordinates
\[
 z_0=qr_0,\qquad z_1=qr_1
\]
for the two lower-level remainders. Pool $z_0,z_1$ together with the five
higher-level remainders under denominator $D'$. The seven-entry budget
forces $|z_i|<D'/2$, hence $|r_i|<D/2$ because $q>0$. The original
denominator congruences therefore recover the same lower-level residues.
Conversely, $q\rho_i<D'/8$ and the other five canonical residues have the
same bound, so the seven-entry budget has witnesses.

This replaces six guard coordinates by three and adds two product
coordinates, saving one variable. Its 260-variable intermediate has
186 rows. The polynomial expressions are still quadratic before taking
their sum of squares.

\section{An odd-square radix and a different recovery order}
\label{sec:odd}

The parent has one initial scalar sign guard on the factorial digit's
output slack $o$, with $q=B+o+1$. Remove that three-coordinate guard and
add the single quadratic row
\begin{equation}
 q=(2z+1)^2,\qquad z\in\Z.
 \label{eq:oddsquare}
\end{equation}
This saves two more variables. It is a restriction on admissible raw
witnesses, so it requires a new completeness argument as well as soundness.

\subsection{Soundness without an initial lower bound of two on the radix}

Initially~\eqref{eq:oddsquare} gives only $q\ge1$. The signed semantic
anchor and square-frame argument orient all principal Pell pairs without
using the deleted sign of $o$. Interpret $V_0=(F^2+1)^t$ first: its
exponent $t\ge16$ is already known, and the base-one case is allowed.
The two retained ceiling signs then give $F\ge1$, $B>F$, and
$y=V_0(V_0-1)$.

Now use the existing request $N=B^x$. Since $B\ge2$, the signed-power
lemma forces $x\ge0$. The first beta endpoint gives $x\ge F+1\ge2$,
and the last gives $x\ge2+(t+1)y$. The numerator requests have base
$q+1\ge2$, so their exponents are nonnegative actual representatives;
the anchor index is nonzero. Consequently $H_0\ge1$. The request with
base $BH_0\ge2$ then gives
\[
 q=(BH_0)^x\ge B^2>B.
\]
Thus $o=q-B-1\ge0$ is recovered. From this point, all previous power
and digit arguments apply with $q\ge2$. The order is
\[
 q\ge1\ \Rightarrow\ V_0\ \Rightarrow\ B\ge2
 \ \Rightarrow\ x\ge2\ \Rightarrow\ H_0\ge1
 \ \Rightarrow\ q\ge B^2>B.
\]
No step uses the very sign it is meant to establish.

\subsection{Completeness by compatible parity choices}

Begin with any finite refutation and the parent's cofinal construction.
Choose its horizon $t\ge16$, $F=t!$, $V_0=(F^2+1)^t$ and
$y=V_0(V_0-1)$. The number $y$ is even, so every beta CRT modulus
$1+(i+1)y$ is odd. The beta code $x$ can therefore be chosen even and
arbitrarily large in its required residue class. Its endpoint and quotient
conditions are preserved by that choice.

The factors $d_i=1+F^2/i$ are odd: $i\mid F$ and
$F^2/i=F(F/i)$ is even. Therefore $B=\prod_i d_i$ and $N=B^x$ are odd.
The independent representative CRT moduli have odd product $N^2$, so
each representative's parity can be prescribed independently.

Prescribe opposite parities to $a,b$, as required to make the semantic
anchor $h$ even. Use one of the other seven representatives to make the
sum of all nine representatives even, hence $H_0$ odd. Ordinary
representatives remain in $(0,2N^2]$. The dominant representative is
chosen as $a_0+2N^2$ or $a_0+3N^2$, maintaining the required parity and
$2N^2\le a<4N^2$. All inherited dominance and cap estimates remain valid.

Now $BH_0$ is odd and $x$ even, so
\[
 z=\frac{(BH_0)^{x/2}-1}{2}\in\Z
\]
supplies~\eqref{eq:oddsquare}. The shared remainder budgets can be
completed as before. This proves completeness of the 258-variable graph.

\section{Smaller coefficients by reversible residual operations}

The preceding graph is converted into literal quadratic residuals by the
retained Horn factorization and integral monic coordinate elimination.
Every eliminated coordinate has an explicit polynomial restoration.
Integral coordinate translations are bijective over the integer domain.
After those steps the main equation has 258 variables, 186 rows, 5,319
terms and height 128,446.

We then apply 30 elementary transformations to its ordered residual list:
\begin{equation}
 r_i\longleftarrow r_i+k r_j\qquad(i\ne j, k\in\Z).
 \label{eq:rowop}
\end{equation}
Each is a unimodular row transformation, with inverse obtained by
replacing $k$ by $-k$. Therefore it preserves the simultaneous zero set
of the residuals over every commutative ring. Over $\Z$ the sum of their
squares vanishes exactly when each residual vanishes, so the transformed
single equation has exactly the same integer zero set.

The search for useful operations was finite and heuristic. Its output,
not the search heuristic, is the certificate: every operation and its
reverse is supplied and checked exactly. These operations lower the
largest residual coefficient from 256 to 128 and the collected quartic
height to 47,281. The final support count returns to 5,319. No claim of
optimal row basis, optimal integer translation, or globally minimal
coefficient height is made.

This is a useful separation of concerns. Number theory justifies deleting
sign witnesses. Reversible algebra improves the displayed coefficients.
Neither improvement is inferred from testing whether the final equation
has a solution.

\section{A compact companion: deleting the auxiliary shift}
\label{sec:compact-shift}

Two compact companions have ten strictly positive integer unknowns.  They
retain the entire fixed indexed BFS943 coefficient tuple of the incoming
\texttt{ConZFC\_Certificate\_Reductions\_2026-10-25.zip} at the repository
snapshot used in this report.  Their live binary arithmetic costs and exact
degrees are
\[
 (77,630)\qquad\hbox{and}\qquad(78,468).
\]
Each saves one addition at unchanged degree from its matched $(78,630)$ or
$(79,468)$ parent. This is a small arithmetic reduction. The new reasoning
is a low-index exclusion that permits signed transport before the input
coordinates have acquired their intended interpretation. It proves that a
previous shift $T+\alpha$ was redundant.

\subsection{The complete polynomial}

The only independent unknowns are
\[
 J,F,\alpha,i,T,w,s,n,\eta,\zeta\in\mathbb Z_{>0}.
\]
The four fixed integers $\beta,K_0,M_C,M_F$ are specified by the finite
recipe below. Every other symbol in the following display is a polynomial
abbreviation, and adds no unknown:
\begin{gather*}
 q=\beta J+1,\quad h=q-1,\quad C=q-F-\alpha,\quad X=wq,\\
 Y=q^3s,\quad k=\eta+\zeta,\quad c=kY+\eta,\quad D=Y^2,\\
 A=Y(X+1)+1,\quad \Delta=A^2-1,\quad H=4A-5,\\
 V_0=XD,\qquad E=1+V_0(V_0+1)k^2,\\
 R=[h(q-F)+\alpha](q^2-1)+(M_C+qM_F)J,\\
 e=k-R-1,\quad B=(K_0+w)C-F,\quad U=ic^2,\\
 a=\Delta U^2E,\quad F_0=1+a,\quad G=EF_0,\\
 F_*=1+4aF_0,\quad Q=\Delta(F_*-1),\\
 b=cT-R,\quad V=F_*b-c,\quad K_+=X(3c-2X)+2.
\end{gather*}
For the 77-operation equation choose $L=V^2$; for the 78-operation
equation choose $L=bV$. Set
\begin{align*}
 M&=hDH,\qquad R_G=M^2G,\\
 Z&=hD\bigl[H(n-LB)+K_+\bigr]-H(DB+he).
\end{align*}
The single equation is
\begin{equation}
 \boxed{\;P_L=Q(Z^2-R_GV^2)+(R_G-Z^2)=0.\;}
 \label{eq:compact-shift-deleted}
\end{equation}
The circuits use only binary addition, subtraction and multiplication;
squares and multiplication by a fixed integer each count as one operation.
The sole change from the corresponding parent is $b=cT-R$ in place of
$c(T+\alpha)-R$.

\paragraph{The theorem and its proof boundary.}
For either choice of $L$, ordinary arithmetic formalizable in ZFC gives
an explicit bijection of the positive zero sets of the new and parent
polynomials. Composing it with the retained parent calibration gives
\[
 \mathrm{ZFC}\vdash
 \left(\neg\exists\mathbf x\in\mathbb Z_{>0}^{10}\;
 P_L(\mathbf x)=0\right)
 \longleftrightarrow\mathrm{Con}(\mathrm{ZFC}).
\]
The compact branch retains its native compiler and K215 ancestry; it is
separate from the direct proof-calculus construction. This statement is an
ordinary proof with named inherited dependencies, not a newly emitted
end-to-end proof-assistant certificate.

\subsection{Why deleting the shift introduces no spurious positive zeros}

There is an exact free-polynomial identity
\[
 P_{\rm new}(T)=P_{\rm parent}(T-\alpha).
\]
To use it over positive integers we must prove $T>\alpha$ at every new
positive zero. It is not valid to assume this inequality as part of the
coordinate change.

The unchanged product-square argument works for either sign of $V$. It
gives positive integers $g,y,f,I$ and $\varepsilon\in\{-1,1\}$ with
\begin{gather*}
 G=g^2,\quad f=1+2a,\quad I=2Ug,\quad f^2=F_*,\\
 f^2-\Delta I^2=1,\quad Q=(\Delta I)^2,\quad c^2\mid I,\\
 QV^2-(Q-1)y^2=1,\qquad Z/M=\varepsilon gy.
\end{gather*}
The signed all-index folding lemma recovers $c=\psi_A(p)$ with $p$ odd
and positive. Here $\psi_A$ denotes the Pell ordinate for $A^2-1$.
The first norm gives
\[
 k=2\psi_{P_0}(n_0),\qquad P_0=2XY^2+1.
\]
The possibility $n_0=1$ would give $c=2Y+1$ strictly between $1$ and
$2A$, so $n_0\ge2$. Consequently
\[
 c>kY>8XY^3,\qquad c>X,\qquad K_+>0.
\]

\paragraph{A useful low-index exclusion.}
Suppose $p=3$, so $c=4A^2-1$. If $n_0\ge3$, then
$k>32X^2Y^4$, whereas $A<2XY$ gives $c<16X^2Y^2$; this contradicts
$c>kY$. Thus $n_0=2$ and $k=8XY^2+4$. The exact identities
\begin{align*}
 2c-(X+2)k&=8Y^2+16Y(X+1)-4X-2>0,\\
 (X+3)k-2c&=8Y[(X-1)Y-2(X+1)]+4X+6>0
\end{align*}
hold for our $X\ge32,Y\ge32^3$. Together with $Y<c/k<Y+1$ they imply
$0<2Y-X<3$. This is impossible since $q\mid X,Y$ and $q\ge32$.
Hence $p\ge5$ and $c>4\Delta+3$. The same exclusion works for
$X\ge3,Y\ge4$ whenever $X,Y$ share any divisor at least three.

\paragraph{A bound that allows both signs.}
Write $S=\Delta I$. Positive Pell classification gives
$S|V|=\chi_S(j)$ and $y=\psi_S(j)$ for $j\ge1$ odd; even $j$
is impossible modulo $S$. Also $f^2>c+4\Delta+3$.
For $j=1$, the equation $f^2b=c\pm1$ is impossible between consecutive
multiples of $f^2$. For $j=3$, the equation becomes
\[
 f^2(b-4\varepsilon_v\Delta)
     =c-\varepsilon_v(4\Delta+3),\qquad\varepsilon_v\in\{-1,1\}.
\]
Its right side is strictly between $0$ and $f^2$ for both signs, using
$c>4\Delta+3$. Thus $j\ge5$, and with $v=|V|$,
\[
 v\ge16Q^2-20Q+5>8Q^2.
\]
Since $c<f^2<Q$ and $f^2\ge3$,
\begin{align*}
 |b|&\ge(v-c)/f^2>(8Q^2-Q)/f^2\\
 &=\Delta(1-f^{-2})(8Q-1)
   \ge\tfrac23\Delta(8Q-1)>4\Delta Q>4g.
\end{align*}
Furthermore $V$ has the sign of $b$ and $v>|b|$. The auxiliary norm
implies $y<3v/2$. Therefore both allowed weights are positive and satisfy
\[
 gy<L/2.
\]

Put $\kappa=K_+/H>0$, $\rho=e/D$, $t=B/h$, and $W=1+hL$.
The exact carrier identity is
$Z/M=n+\kappa-Wt-\rho$, and
$b-e=cT-k+1>0$ gives $\rho<b/D$ even if $b<0$.
If $B\le-1$, then $Wt\le-L-1/h$, whence
\[
 n=\varepsilon gy-\kappa+Wt+\rho<-L/2+b/D<0.
\]
For $b<0$ the final inequality is immediate; for $b>0$ it follows from
$D>2$ and $L>b$. This contradicts $n>0$. Thus $B\ge0$, so $C>0$ and
$\alpha<q$. The fixed mask bounds now give $0<R<q^4$, while
$c>8XY^3\ge8q^{10}$. Consequently $b=cT-R>c/2>0$.
Finally
\[
 g^2=E+\Delta U^2E^2>(A-1)^2U^2E^2,
 \qquad g/c>(A-1)c\ge q,
\]
and therefore
\[
 T=(b+R)/c>b/c>4g/c>q>\alpha.
\]
The positive inverse coordinate $T_{\rm parent}=T-\alpha$ is now
justified. In the other direction, every positive parent zero maps to a
new positive zero by $T_{\rm new}=T_{\rm parent}+\alpha$. These maps
are inverse bijections of the whole positive zero sets.

\paragraph{Exact arithmetic checks.}
The entire leading homogeneous components, with $q_0=\beta J$ and
$k_0=\eta+\zeta$, are
\begin{gather*}
 2^{14}q_0^{382}w^{46}s^{112}i^{20}k_0^{64}T^4
       (q_0-F-\alpha)^2,\\
 -2^{10}q_0^{288}w^{34}s^{84}i^{14}k_0^{46}T^2.
\end{gather*}
They are nonzero and have degrees $630$ and $468$. The verifier checks
all actual circuit registers through 28 free-polynomial identities per
equation, the complete leading forms, the fixed coefficient reconstruction,
the three new low-index identities, and ten changed-formula rejection
controls per equation. The full proof and exact dependency references are
in \texttt{compact/PROOF.md}; the incoming ordinary proof and explicit
signed folding lemma are included. Finite artifact checks do not prove the
unbounded statements by sampling.

\begingroup
\newcommand{\cfl}[1]{\left\lfloor#1\right\rfloor}
\newcommand{\crange}[2]{\{#1,\ldots,#2\}}
\subsection{The complete finite coefficient recipe}
Every index, set, sum, product and floor below defines a fixed constant, not
another unknown. The letter $B_\ast$ belongs only to this constant recipe;
it is distinct from the polynomial abbreviation $B$ above. Put
\[
\rho=2^{67},\quad B_\ast=\rho^{8807},\quad\beta=B_\ast-1,\quad
G_m(x)=\sum_{j=0}^{m-1}x^j.
\]
Define
\[
\begin{split}
\mathcal E={}&\crange0{942}\cup\crange{983}{992}\cup\crange{1003}{1011}\\
&\cup\crange{1013}{1031}\cup\{1063,1064,1095,2195,3296\}.
\end{split}
\]
Then
\begin{align*}
M_C={}&B_\ast-1-\sum_{e\in\mathcal E\setminus\{1\}}\rho^e-2\rho^{1063},\\
M_F={}&B_\ast-1+4+(2^{51}-2)(\rho^{4393}+\rho^{5495})\\
&+\rho^{1013}G_{10}(\rho)+\rho^{1065}G_9(\rho)
  +\rho^{2195}+\rho^{3296}\\
&+(2^{65}-1)\rho^{8792}G_{13}(\rho)+(2^{19}-1)\rho^{8805},\\
K_0={}&\rho^{30}+\rho^{62}+\rho^{1100}+\rho^{2201}+\rho^{8849}\\
&+(\rho^{4393}+\rho^{5495})\Phi(1/\rho),
\end{align*}
where
\[
\begin{split}
\Phi(x)={}&(1+2^{51})G_{943}(x)
 -2^{10}\sum_{i=0}^{942}v(t_i)x^i\\
&+\sum_{(b,e,m)\in\mathcal C}2^b x^e G_m(2x),\\
\mathcal C={}&\{(10,1023,8),(18,1013,9),(27,1003,8),\\
&(35,983,9),(44,1031,1),(45,1022,1),(46,1011,1),\\
&(47,992,1),(48,1095,1),(49,2195,1),(50,3296,1)\}.
\end{split}
\]
The rational intermediate in $K_0$ cancels exactly; all four results are positive integers.

\textbf{The 943 indexed fixed tiles.}
Use the actions $u_j,v_j,d_j$ decoded below. Let
\[
\begin{gathered}
h(q,b)=2(q+1)+b,\quad
\mathcal R=\{u_j:d_j=1,\ 0\le j<430\},\\
\mathcal L=\{u_j:d_j=0,\ 0\le j<430\},\\
r_q=6+|\mathcal R\cap\crange0{q-1}|,\quad
l_q=205+|\mathcal L\cap\crange0{q-1}|.
\end{gathered}
\]
An interval ending at $-1$ is empty. The table defines $t_i$ for
$0\le i\le942$; all index ranges are fixed:
\[
\setlength{\arraycolsep}{1.4pt}\begin{array}{c|l}
\text{index }i&\text{tile }t_i=(w,n,e,s)\\\hline
0,1&(2,435,3,1),\ (5,0,0,0)\\
j+2\ (j=2q+b<430)&((1-d_j)l_{u_j},h(q,b),d_jr_{u_j},v_j)\\
432+b\ (b=0,1)&(0,h(215,b),0,h(215,b))\\
434+b\ (b=0,1)&(0,b,0,b)\\
424+2r_q+b\ (q\in\mathcal R,b=0,1)&(r_q,b,0,h(q,b))\\
424+2l_q+b\ (q\in\mathcal L,b=0,1)&(0,b,l_q,h(q,b))\\
924+j\ (0\le j<4)&(0,j+430\cfl{j/2},0,438)\\
928&(2,434,2,0)\\
929&(3,436,4,4)\\
930&(4,436,4,0)\\
497+t\ (t=434,435,436)&(1,438,1,t)\\
\substack{934+3a+b-6\cfl{a/4}\\a\in\{0,1,4\},\ b=0,1,2}&(a,437,b,437)
\end{array}
\]
For $t=(w,n,e,s)$ put
\[
v(t)=w+2^8n+2^{17}e+2^{25}s+127\,2^{34}.
\]

\textbf{A closed decoder for the 430 actions.}
Put $D_b(A,k)=\cfl{A/b^k}\bmod b$, where
$a\bmod b=a-b\cfl{a/b}$. For $0\le q<215$, $b=0,1$, $j=2q+b$:
\[
\begin{gathered}
N=\mathcal A\bmod2^{636},\quad
T_0=\cfl{\mathcal A/2^{636}}\bmod2^{1675},\\
E_0=\cfl{\mathcal A/2^{2311}},\\
\nu_j=D_2(N,j),\quad \sigma_j=\sum_{k=0}^{j-1}\nu_k,\\
m_q=D_2(N,430+q),\quad \tau_q=\sum_{k=0}^{q-1}m_k,\\
F_q=14(1-m_q)+m_qD_{16}(E_0,\tau_q),\\
u_j=\nu_j(1+\sigma_j)+(1-\nu_j)D_{216}(T_0,j-\sigma_j),\\
v_j=D_2(F_q,2b),\quad d_j=D_2(F_q,2b+1).
\end{gathered}
\]
Empty sums are zero. Here $(v_0,d_0,u_0)=(1,1,1)$; the constructor tile
uses this actual first target, including its state renaming.

\textbf{The integer $\mathcal A$.}
The alphabet, in order of digit values $0,\ldots,61$, is:\par
{\fontsize{7.7}{8.2}\selectfont\ttfamily 0123456789ABCDEFGHIJKLMNOPQRSTUVWXYZabcdefghijklmnopqrstuvwxyz}.\par
Concatenate the following 436 characters as digits $e_0,\ldots,e_{435}$;
put $\mathcal A=\sum_{j=0}^{435}e_j62^{435-j}$.\par
\begingroup\fontsize{7.7}{8.2}\selectfont\ttfamily
BxEnhvy6NipzkWmROo0GsA3uK3twWLxscLq20UTumy5Cdu461LwD6g7RDNzF8Y\par
iYDFTpB8ehAypd7d6OsOL4APamnJeNXQypUezCmW9n4oXP6XfgmiWuAWZ8ZJ7U\par
OP7gu4kVdlGZqZHULQMar1dXNeJW4i7oNVI8usp8BlovjjZbwPCovZI328fRUq\par
4q92R81s5GjWQaQfGhkq9UYMK3rRMaBmsOIBkJqc2cj2VRIZ6mjzkBIOkW9SjZ\par
FzeA6J5u4j3sBL3SdgwXzM9mHvQivu8EJ1TBTpRhHDqZIbmYgwQ8VqTAFotl6F\par
We1YC554GKrVHsA8oUGKSA6F6k5KREBSD0Tgd7bFALY84LLbpw574BQyvdoHrl\par
0Oe35uwDf15DpT6VJWZBjSbHPpPAINYTte0eUxnTFTKBlGj8Am6sb1gLldh8WZ\par
8P\par
\endgroup
\endgroup

The recipe reproduces all four retained integer payloads byte for byte. The
largest fixed input coefficient has 592,884 bits. This is a stored-port
size, not the coefficient height of the fully expanded polynomial.


\section{Reproduction and verification boundaries}

The following files specify the result without needing an external
repository or a symbolic-mathematics service:
\begin{itemize}
\item \code{equations/P258/equation.json}: every ordered variable and all
186 literal residuals; its adjacent expansion gives every quartic term.
\item \code{equations/P261/}: the complete 261-variable companion.
\item \code{direct/}: the literal guarded source, exact changes, source
restorations, translations, arithmetic proof and finite verifier.
\item \code{compact/}: both complete compact circuits, the four fixed
integer payloads and their finite reconstruction, the new proof and
the retained parent proofs.
\item \code{sources/}: pinned inherited arithmetic and logical interfaces.
\item \code{verify.py}: the combined finite check, including reverse row
operations, complete SOS expansion, and parsing the printed residuals.
\end{itemize}

From the extracted archive root, use Python 3.10 or later:
\begin{verbatim}
python3 -B verify.py
python3 -B -O verify.py
python3 -B build_endpoints.py --check
python3 -B evaluate.py point.json
\end{verbatim}
The default evaluator expects exactly 258 signed integer coordinates.
It is an exact evaluator, not a consistency test or a solver. The
standalone LaTeX paper is rebuilt with \code{sh build\_report.sh}.

The mathematical claims have two kinds of support. The unbounded
arguments prove soundness and completeness, including the delicate order
of sign recovery. Finite checks authenticate the entire written
polynomial, transformations, degrees, coefficient heights and circuit
counts. The latter do not prove an unbounded number-theoretic lemma by
sampling. Independent reviews in this packet are separate reviews within
this AI-assisted workflow, not outside human peer review.

Formalization is a natural next step. Existing formalized DPRM work shows
how arithmetic graphs and proof encodings can be handled by a proof
assistant~\cite{dprmformal}. We do not substitute that general theorem
for the explicit coefficients or the calibrated equivalence in this
particular construction.

\section{What is new, and further research questions}

The direct contribution is the completed application of symmetric
remainder bounds to the signed C222 graph, including a pre-factorial radix
margin, a two-scale merger, and the odd-square recovery order. The compact
contribution is the proof that the shift $T+\alpha$ can be removed while
preserving the entire positive zero set by an explicit triangular
bijection. The row operations give a fully certified coefficient reduction
at the new direct endpoints. CRT, Pell equations, ternary nonnegativity
representations and elementary row operations are classical ingredients;
no historical priority for them is asserted.

Several questions remain particularly concrete:
\begin{enumerate}
\item Can all digit levels share one quadratic budget, with fewer added
scaling coordinates than the guard coordinates it removes? The adjacent
levels work because their ratio is the already present radix $q$.
\item Can another cofinal arithmetic choice remove one of the two
remaining ceiling sign graphs? A proposed restriction must preserve
completeness; a convenient square condition cannot simply be assumed.
\item Can the set-theoretic seed and logical vacuity parts of the
numerical calculus be compiled separately to reduce the arithmetic front?
The repository already contains warnings against dropping necessary
vacuity instances. A new front needs its own complete calibration and
explicit finite polynomial, not merely a suggestive formula.
\item Can the signed residue-recognition lemma and the revised recovery
order be formalized first as reusable Lean/Rocq lemmas, then composed
with a verified literal compiler? This would close the principal gap
between the current ordinary proof and an end-to-end formal endpoint.
\item Can coefficient-aware row and coordinate transformations lower
height further without increasing variables or support? The present
30-operation certificate gives a reproducible starting point, not a
lower bound.
\end{enumerate}

\begin{thebibliography}{9}
\bibitem{repo}
V. Reshetnikov and contributors,
\emph{ProveIt}, immutable comparison snapshot
\code{e0c98c629bf10d30b695884f6b1d2ab361bdf513}.
\url{https://github.com/VladimirReshetnikov/ProveIt/tree/e0c98c629bf10d30b695884f6b1d2ab361bdf513}.
\bibitem{cofinal}
\emph{Cofinal anchors and height one equations for ZFC consistency},
round 17, 25 October 2026, incoming research archive in~\cite{repo}.
Its direct ordinary proof and exact guarded source are retained here.
\bibitem{cert}
\emph{Eliminating auxiliary witnesses in Diophantine consistency},
25 October 2026, \code{ConZFC\_Certificate\_Reductions\_2026-10-25.zip}
in~\cite{repo}. The compact component is retained byte for byte.
\bibitem{digits}
\emph{Pooled low digits and a semantic Pell cap},
\code{conzfc\_graph\_digits\_2026-10-24.zip} in~\cite{repo}.
The natural-domain pooling argument is a predecessor; the present signed
remainder recognition and recovery schedule are different.
\bibitem{bayer}
J. Bayer, M. David, M. Ha\ss ler, Y. Matiyasevich and D. Schleicher,
\emph{Diophantine equations over $\Z$: universal bounds and parallel
formalization}, arXiv:2506.20909v3, 22 October 2026.
\url{https://arxiv.org/abs/2506.20909}.
\bibitem{dprmformal}
J. Bayer, M. David, B. Stock, A. Pal, Y. Matiyasevich and D. Schleicher,
\emph{Diophantine Equations and the DPRM Theorem}, Archive of Formal
Proofs, 2022. \url{https://isa-afp.org/entries/DPRM_Theorem.html}.
\end{thebibliography}

\clearpage
\appendix
\section{Every residual of the main equation}\label{app:residuals}

The independent unknowns are exactly $x_0,\ldots,x_{257}\in\Z$.
The equation is $\sum_{j=0}^{185}r_j^2=0$, with the following complete
list. All constants below are literal integers; no reference to a proof
predicate, variable exponent, recursion, or external numerical port is
needed to evaluate the polynomial. The file
\code{equations/P258/variables.tsv} also records the compiler's semantic
names, which are not needed for this definition.

\begingroup\small
\setlength{\abovedisplayskip}{3pt}
\setlength{\belowdisplayskip}{3pt}
\setlength{\jot}{1pt}
% Generated literal quadratic residuals. Parsed back by verify.py.
\[\begin{aligned}
r_{0}&=-x_{3}^{2}+x_{4}
\end{aligned}\]
\[\begin{aligned}
r_{1}&=x_{0}-x_{0}x_{2}-x_{2}+x_{9}
\end{aligned}\]
\[\begin{aligned}
r_{2}&=-1+x_{1}-x_{2}x_{7}-x_{3}
\end{aligned}\]
\[\begin{aligned}
r_{3}&=-1+x_{1}-x_{8}x_{9}
\end{aligned}\]
\[\begin{aligned}
r_{4}&=1+x_{4}-x_{5}x_{20}+x_{18}
\end{aligned}\]
\[\begin{aligned}
r_{5}&=-1+2x_{10}-x_{10}^{2}-2x_{10}x_{11}+2x_{11}-x_{11}^{2}+x_{21}
\end{aligned}\]
\[\begin{aligned}
r_{6}&=-4-12x_{10}-16x_{11}-4x_{21}+x_{22}-x_{37}
\end{aligned}\]
\[\begin{aligned}
r_{7}&=-2+2x_{10}-x_{10}^{2}-4x_{10}x_{22}+4x_{22}-4x_{22}^{2}+x_{23}
\end{aligned}\]
\[\begin{aligned}
r_{8}&=67-60x_{10}-3x_{10}^{2}-60x_{10}x_{22}-24x_{10}x_{23}-116x_{22}-124x_{22}^{2}\\
&{}-64x_{22}x_{23}-31x_{23}-4x_{23}^{2}+x_{24}
\end{aligned}\]
\[\begin{aligned}
r_{9}&=-2-x_{11}^{2}-2x_{11}x_{12}-x_{12}^{2}+x_{25}
\end{aligned}\]
\[\begin{aligned}
r_{10}&=-9+7x_{10}-x_{10}^{2}-4x_{10}x_{11}-4x_{10}x_{25}+12x_{11}-4x_{11}^{2}-8x_{11}x_{25}+12x_{25}\\
&{}-4x_{25}^{2}+x_{26}
\end{aligned}\]
\[\begin{aligned}
r_{11}&=-1+2x_{10}-x_{10}^{2}-2x_{10}x_{12}+2x_{12}-x_{12}^{2}+x_{27}
\end{aligned}\]
\[\begin{aligned}
r_{12}&=2-26x_{10}-68x_{10}^{2}-24x_{10}x_{11}-64x_{10}x_{12}+48x_{10}x_{27}+48x_{10}x_{30}\\
&{}+52x_{11}^{2}+32x_{11}x_{21}+96x_{11}x_{27}+64x_{11}x_{30}-48x_{12}-16x_{12}^{2}+32x_{12}x_{27}\\
&{}+32x_{12}x_{30}+6x_{21}+4x_{21}^{2}+24x_{21}x_{27}+16x_{21}x_{30}-32x_{27}+28x_{27}^{2}\\
&{}+32x_{27}x_{30}+x_{28}-8x_{30}+8x_{30}^{2}-2x_{31}
\end{aligned}\]
\[\begin{aligned}
r_{13}&=-4x_{26}^{2}-8x_{26}x_{28}-4x_{28}^{2}+x_{29}
\end{aligned}\]
\[\begin{aligned}
r_{14}&=-6+16x_{12}-16x_{12}^{2}-16x_{12}x_{25}-16x_{12}x_{27}+8x_{25}-4x_{25}^{2}-8x_{25}x_{27}\\
&{}+8x_{27}-4x_{27}^{2}+x_{30}
\end{aligned}\]
\[\begin{aligned}
r_{15}&=27-46x_{10}-10x_{10}^{2}-54x_{10}x_{11}-38x_{10}x_{12}-24x_{10}x_{21}-48x_{10}x_{27}\\
&{}-24x_{10}x_{30}-34x_{11}-43x_{11}^{2}-64x_{11}x_{12}-32x_{11}x_{21}-64x_{11}x_{27}-32x_{11}x_{30}\\
&{}-2x_{12}-11x_{12}^{2}-16x_{12}x_{21}-32x_{12}x_{27}-16x_{12}x_{30}-19x_{21}-4x_{21}^{2}\\
&{}-16x_{21}x_{27}-8x_{21}x_{30}+3x_{27}-16x_{27}^{2}-16x_{27}x_{30}+4x_{30}-4x_{30}^{2}+x_{31}
\end{aligned}\]
\[\begin{aligned}
r_{16}&=46-5x_{10}-x_{10}^{2}-12x_{10}x_{11}-16x_{10}x_{12}-4x_{10}x_{25}-36x_{11}-8x_{11}^{2}\\
&{}-40x_{11}x_{12}-24x_{11}x_{25}-48x_{12}-36x_{12}^{2}-32x_{12}x_{25}-40x_{25}-4x_{25}^{2}+x_{32}
\end{aligned}\]
\[\begin{aligned}
r_{17}&=24-44x_{10}-26x_{10}^{2}-12x_{10}x_{11}-76x_{10}x_{12}-24x_{10}x_{27}-3x_{11}+7x_{11}^{2}\\
&{}-4x_{11}x_{27}-52x_{12}-46x_{12}^{2}-32x_{12}x_{27}-8x_{25}-18x_{27}-4x_{27}^{2}+x_{33}
\end{aligned}\]
\[\begin{aligned}
r_{18}&=2-4x_{32}^{2}-8x_{32}x_{33}-4x_{33}^{2}+x_{34}
\end{aligned}\]
\[\begin{aligned}
r_{19}&=-12+16x_{10}-4x_{10}^{2}-8x_{10}x_{12}-8x_{10}x_{21}-8x_{10}x_{25}+16x_{12}-4x_{12}^{2}\\
&{}-8x_{12}x_{21}-8x_{12}x_{25}+16x_{21}-4x_{21}^{2}-8x_{21}x_{25}+16x_{25}-4x_{25}^{2}+x_{35}
\end{aligned}\]
\[\begin{aligned}
r_{20}&=22-74x_{10}-39x_{10}^{2}-110x_{10}x_{12}-64x_{10}x_{25}-32x_{10}x_{27}-32x_{10}x_{35}\\
&{}-76x_{12}-75x_{12}^{2}-80x_{12}x_{25}-40x_{12}x_{27}-40x_{12}x_{35}-8x_{25}-16x_{25}^{2}\\
&{}-16x_{25}x_{27}-16x_{25}x_{35}-31x_{27}-4x_{27}^{2}-8x_{27}x_{35}-4x_{35}-4x_{35}^{2}+x_{36}
\end{aligned}\]
\[\begin{aligned}
r_{21}&=25-42x_{10}-24x_{10}^{2}-72x_{10}x_{11}-24x_{10}x_{21}-48x_{11}-45x_{11}^{2}+14x_{11}x_{12}\\
&{}-32x_{11}x_{21}+7x_{12}^{2}-18x_{21}-4x_{21}^{2}-7x_{25}+x_{37}
\end{aligned}\]
\[\begin{aligned}
r_{22}&=-2+2x_{10}-x_{10}^{2}-4x_{10}x_{37}+4x_{37}-4x_{37}^{2}+x_{38}
\end{aligned}\]
\[\begin{aligned}
r_{23}&=67-60x_{10}-3x_{10}^{2}-60x_{10}x_{37}-24x_{10}x_{38}-116x_{37}-124x_{37}^{2}\\
&{}-64x_{37}x_{38}-31x_{38}-4x_{38}^{2}+x_{39}
\end{aligned}\]
\[\begin{aligned}
r_{24}&=-x_{11}^{2}-4x_{11}x_{39}-4x_{39}^{2}+x_{40}
\end{aligned}\]
\[\begin{aligned}
r_{25}&=10-30x_{10}-64x_{10}^{2}-128x_{10}x_{11}-64x_{10}x_{21}-32x_{11}-56x_{11}^{2}\\
&{}+16x_{11}x_{12}-64x_{11}x_{21}+8x_{12}^{2}-14x_{21}-16x_{21}^{2}-8x_{25}+x_{41}
\end{aligned}\]
\[\begin{aligned}
r_{26}&=52+8x_{11}-36x_{11}^{2}-40x_{11}x_{12}-32x_{11}x_{25}-32x_{11}x_{41}+6x_{12}-8x_{12}^{2}\\
&{}-24x_{12}x_{25}-24x_{12}x_{41}-26x_{25}-4x_{25}^{2}-8x_{25}x_{41}+4x_{41}-4x_{41}^{2}+x_{42}
\end{aligned}\]
\[\begin{aligned}
r_{27}&=8x_{10}-64x_{10}^{2}-96x_{10}x_{12}-32x_{10}x_{27}-32x_{10}x_{42}+6x_{12}-36x_{12}^{2}\\
&{}-24x_{12}x_{27}-24x_{12}x_{42}+2x_{27}-4x_{27}^{2}-8x_{27}x_{42}+4x_{42}-4x_{42}^{2}+x_{43}
\end{aligned}\]
\[\begin{aligned}
r_{28}&=-x_{12}^{2}-2x_{12}x_{13}-x_{13}^{2}+x_{44}
\end{aligned}\]
\[\begin{aligned}
r_{29}&=2x_{10}-x_{10}^{2}-4x_{10}x_{12}-4x_{10}x_{44}+4x_{12}-4x_{12}^{2}-8x_{12}x_{44}+4x_{44}\\
&{}-4x_{44}^{2}+x_{45}
\end{aligned}\]
\[\begin{aligned}
r_{30}&=-26+17x_{10}-3x_{10}^{2}-4x_{10}x_{11}-4x_{10}x_{25}-8x_{10}x_{27}+12x_{11}-8x_{11}x_{27}\\
&{}+12x_{25}-8x_{25}x_{27}-x_{26}+24x_{27}-4x_{27}^{2}+x_{46}
\end{aligned}\]
\[\begin{aligned}
r_{31}&=-8x_{10}-12x_{10}^{2}-16x_{10}x_{12}-32x_{10}x_{27}-16x_{10}x_{44}-16x_{10}x_{45}-16x_{10}x_{46}\\
&{}-20x_{12}-32x_{12}x_{27}-16x_{12}x_{45}-16x_{12}x_{46}-16x_{27}^{2}-32x_{27}x_{44}-16x_{27}x_{45}\\
&{}-16x_{27}x_{46}-20x_{44}-16x_{44}x_{45}-16x_{44}x_{46}-6x_{45}-4x_{45}^{2}-8x_{45}x_{46}\\
&{}-4x_{46}^{2}+x_{47}
\end{aligned}\]
\[\begin{aligned}
r_{32}&=x_{12}-25x_{12}^{2}-40x_{12}x_{27}-20x_{12}x_{30}-16x_{27}^{2}-16x_{27}x_{30}-4x_{30}^{2}\\
&{}+x_{48}
\end{aligned}\]
\[\begin{aligned}
r_{33}&=-1+8x_{11}-4x_{11}^{2}-8x_{11}x_{25}+8x_{25}-4x_{25}^{2}+x_{49}
\end{aligned}\]
\[\begin{aligned}
r_{34}&=-1-16x_{11}^{2}-8x_{11}x_{12}-32x_{11}x_{25}-16x_{11}x_{49}-x_{12}^{2}-8x_{12}x_{25}-4x_{12}x_{49}\\
&{}-16x_{25}^{2}-16x_{25}x_{49}-4x_{49}^{2}+x_{50}
\end{aligned}\]
\[\begin{aligned}
r_{35}&=1+x_{11}-x_{11}^{2}-4x_{11}x_{12}-4x_{11}x_{50}-4x_{12}^{2}-8x_{12}x_{50}-4x_{50}^{2}+x_{51}
\end{aligned}\]
\[\begin{aligned}
r_{36}&=-1-4x_{51}-4x_{51}^{2}+x_{52}
\end{aligned}\]
\[\begin{aligned}
r_{37}&=-1-x_{12}^{2}-8x_{12}x_{51}-4x_{12}x_{52}-16x_{51}^{2}-16x_{51}x_{52}-4x_{52}^{2}+x_{53}
\end{aligned}\]
\[\begin{aligned}
r_{38}&=2+x_{11}-x_{11}^{2}-4x_{11}x_{12}-4x_{11}x_{53}-4x_{12}^{2}-8x_{12}x_{53}-4x_{53}^{2}+x_{54}
\end{aligned}\]
\[\begin{aligned}
r_{39}&=-4x_{48}^{2}-8x_{48}x_{54}-4x_{54}^{2}+x_{55}
\end{aligned}\]
\[\begin{aligned}
r_{40}&=-9+7x_{10}-x_{10}^{2}-8x_{10}x_{48}-4x_{10}x_{55}+24x_{48}-16x_{48}^{2}-16x_{48}x_{55}\\
&{}+12x_{55}-4x_{55}^{2}+x_{56}
\end{aligned}\]
\[\begin{aligned}
r_{41}&=-1-4x_{56}-4x_{56}^{2}+x_{57}
\end{aligned}\]
\[\begin{aligned}
r_{42}&=-x_{12}^{2}-8x_{12}x_{56}-4x_{12}x_{57}-16x_{56}^{2}-16x_{56}x_{57}-4x_{57}^{2}+x_{58}
\end{aligned}\]
\[\begin{aligned}
r_{43}&=-x_{10}-x_{10}^{2}-4x_{10}x_{12}-4x_{10}x_{58}-4x_{12}-4x_{12}^{2}-8x_{12}x_{58}-4x_{58}\\
&{}-4x_{58}^{2}+x_{59}
\end{aligned}\]
\[\begin{aligned}
r_{44}&=-10-12x_{59}-4x_{59}^{2}+x_{60}
\end{aligned}\]
\[\begin{aligned}
r_{45}&=1+4x_{11}-4x_{11}^{2}-8x_{11}x_{25}-16x_{11}x_{59}-8x_{11}x_{60}+4x_{25}-4x_{25}^{2}\\
&{}-16x_{25}x_{59}-8x_{25}x_{60}+8x_{59}-16x_{59}^{2}-16x_{59}x_{60}+4x_{60}-4x_{60}^{2}+x_{61}
\end{aligned}\]
\[\begin{aligned}
r_{46}&=x_{11}-25x_{11}^{2}-40x_{11}x_{25}-20x_{11}x_{61}-16x_{25}^{2}-16x_{25}x_{61}-4x_{61}^{2}\\
&{}+x_{62}
\end{aligned}\]
\[\begin{aligned}
r_{47}&=-10-12x_{54}-4x_{54}^{2}+x_{63}
\end{aligned}\]
\[\begin{aligned}
r_{48}&=1-16x_{54}^{2}-16x_{54}x_{62}-16x_{54}x_{63}-4x_{62}^{2}-8x_{62}x_{63}-4x_{63}^{2}+x_{64}
\end{aligned}\]
\[\begin{aligned}
r_{49}&=x_{12}-x_{12}^{2}-8x_{12}x_{62}-4x_{12}x_{64}-16x_{62}^{2}-16x_{62}x_{64}-4x_{64}^{2}+x_{65}
\end{aligned}\]
\[\begin{aligned}
r_{50}&=-4x_{10}-4x_{10}^{2}-4x_{10}x_{12}-8x_{10}x_{21}-x_{12}-x_{12}^{2}-4x_{12}x_{21}-4x_{21}\\
&{}-4x_{21}^{2}+x_{66}
\end{aligned}\]
\[\begin{aligned}
r_{51}&=1+4x_{10}-4x_{10}^{2}-8x_{10}x_{21}-8x_{10}x_{66}+4x_{21}-4x_{21}^{2}-8x_{21}x_{66}+4x_{66}\\
&{}-4x_{66}^{2}+x_{67}
\end{aligned}\]
\[\begin{aligned}
r_{52}&=-2-4x_{10}-16x_{10}^{2}-32x_{10}x_{21}-16x_{10}x_{65}-16x_{10}x_{67}-4x_{21}-16x_{21}^{2}\\
&{}-16x_{21}x_{65}-16x_{21}x_{67}-4x_{65}^{2}-8x_{65}x_{67}-2x_{67}-4x_{67}^{2}+x_{68}
\end{aligned}\]
\[\begin{aligned}
r_{53}&=-2-4x_{68}-4x_{68}^{2}+x_{69}
\end{aligned}\]
\[\begin{aligned}
r_{54}&=-x_{3}x_{15}+x_{70}
\end{aligned}\]
\[\begin{aligned}
r_{55}&=x_{14}-x_{15}-2x_{24}x_{70}-x_{71}
\end{aligned}\]
\[\begin{aligned}
r_{56}&=x_{14}-x_{15}-4x_{26}x_{70}-2x_{29}x_{70}-x_{72}
\end{aligned}\]
\[\begin{aligned}
r_{57}&=x_{14}-x_{15}-2x_{31}x_{70}-6x_{70}-x_{73}
\end{aligned}\]
\[\begin{aligned}
r_{58}&=x_{14}-x_{15}-4x_{32}x_{70}-2x_{34}x_{70}-x_{74}
\end{aligned}\]
\[\begin{aligned}
r_{59}&=x_{14}-x_{15}-2x_{36}x_{70}+6x_{70}-x_{75}
\end{aligned}\]
\[\begin{aligned}
r_{60}&=x_{14}-x_{15}-4x_{39}x_{70}-2x_{40}x_{70}-x_{76}
\end{aligned}\]
\[\begin{aligned}
r_{61}&=x_{14}-x_{15}-2x_{43}x_{70}-x_{77}
\end{aligned}\]
\[\begin{aligned}
r_{62}&=x_{14}-x_{15}-2x_{47}x_{70}-4x_{70}-x_{78}
\end{aligned}\]
\[\begin{aligned}
r_{63}&=-x_{10}x_{11}+x_{10}x_{12}+x_{11}+x_{11}x_{12}-x_{12}-x_{12}^{2}+x_{79}
\end{aligned}\]
\[\begin{aligned}
r_{64}&=1+2x_{13}+x_{13}^{2}-x_{79}^{2}+x_{80}
\end{aligned}\]
\[\begin{aligned}
r_{65}&=-x_{80}^{2}+x_{81}
\end{aligned}\]
\[\begin{aligned}
r_{66}&=-4x_{68}x_{70}-2x_{69}x_{70}+2x_{70}+x_{82}
\end{aligned}\]
\[\begin{aligned}
r_{67}&=-x_{14}^{2}+2x_{14}x_{15}+2x_{14}x_{82}-x_{15}^{2}-2x_{15}x_{82}-x_{82}^{2}+x_{83}
\end{aligned}\]
\[\begin{aligned}
r_{68}&=-1-x_{3}x_{11}+x_{84}
\end{aligned}\]
\[\begin{aligned}
r_{69}&=-x_{15}x_{84}+x_{86}x_{87}+x_{88}
\end{aligned}\]
\[\begin{aligned}
r_{70}&=-1+x_{3}-x_{3}x_{10}+x_{85}
\end{aligned}\]
\[\begin{aligned}
r_{71}&=-1+2x_{3}-2x_{3}x_{10}-2x_{3}x_{21}+x_{86}
\end{aligned}\]
\[\begin{aligned}
r_{72}&=-x_{14}x_{85}+x_{87}
\end{aligned}\]
\[\begin{aligned}
r_{73}&=-1-6x_{3}x_{10}-8x_{3}x_{11}-2x_{3}x_{21}+x_{89}
\end{aligned}\]
\[\begin{aligned}
r_{74}&=x_{14}x_{84}-x_{15}x_{89}+x_{90}
\end{aligned}\]
\[\begin{aligned}
r_{75}&=-x_{14}x_{71}+x_{15}x_{71}+x_{91}
\end{aligned}\]
\[\begin{aligned}
r_{76}&=-x_{72}x_{91}+x_{92}
\end{aligned}\]
\[\begin{aligned}
r_{77}&=-x_{73}x_{92}+x_{93}
\end{aligned}\]
\[\begin{aligned}
r_{78}&=-x_{74}x_{93}+x_{94}
\end{aligned}\]
\[\begin{aligned}
r_{79}&=-x_{75}x_{94}+x_{95}
\end{aligned}\]
\[\begin{aligned}
r_{80}&=-x_{76}x_{95}+x_{96}
\end{aligned}\]
\[\begin{aligned}
r_{81}&=-x_{77}x_{96}+x_{97}
\end{aligned}\]
\[\begin{aligned}
r_{82}&=-x_{78}x_{97}+x_{98}
\end{aligned}\]
\[\begin{aligned}
r_{83}&=-x_{88}x_{98}+x_{99}
\end{aligned}\]
\[\begin{aligned}
r_{84}&=-x_{90}x_{99}+x_{100}
\end{aligned}\]
\[\begin{aligned}
r_{85}&=-x_{100}^{2}+x_{101}
\end{aligned}\]
\[\begin{aligned}
r_{86}&=-x_{2}-x_{2}x_{18}+x_{18}+x_{102}
\end{aligned}\]
\[\begin{aligned}
r_{87}&=-x_{16}x_{102}+x_{103}
\end{aligned}\]
\[\begin{aligned}
r_{88}&=-2x_{1}^{2}+2x_{1}x_{14}+2x_{1}x_{15}+2x_{1}x_{103}+2x_{1}x_{104}-x_{14}^{2}-2x_{14}x_{103}\\
&{}-x_{15}^{2}-2x_{15}x_{104}+x_{19}x_{106}-x_{81}x_{101}-x_{83}x_{101}-x_{103}^{2}-x_{104}^{2}
\end{aligned}\]
\[\begin{aligned}
r_{89}&=-x_{2}x_{17}+x_{17}-x_{17}x_{102}+x_{104}
\end{aligned}\]
\[\begin{aligned}
r_{90}&=-9x_{10}-9x_{11}-4x_{21}-2x_{24}-4x_{26}-2x_{29}-2x_{31}-4x_{32}-2x_{34}-2x_{36}-4x_{39}\\
&{}-2x_{40}-2x_{43}-2x_{47}-4x_{68}-2x_{69}+x_{105}
\end{aligned}\]
\[\begin{aligned}
r_{91}&=-x_{107}x_{110}+x_{110}+x_{112}
\end{aligned}\]
\[\begin{aligned}
r_{92}&=-x_{5}x_{109}+x_{108}-x_{109}x_{112}-x_{111}
\end{aligned}\]
\[\begin{aligned}
r_{93}&=-x_{19}x_{114}+x_{117}
\end{aligned}\]
\[\begin{aligned}
r_{94}&=x_{113}-x_{115}x_{117}-x_{116}
\end{aligned}\]
\[\begin{aligned}
r_{95}&=-x_{117}x_{119}+x_{118}-x_{120}
\end{aligned}\]
\[\begin{aligned}
r_{96}&=-x_{117}x_{122}+x_{121}-x_{123}
\end{aligned}\]
\[\begin{aligned}
r_{97}&=-x_{117}x_{125}+x_{124}-x_{126}
\end{aligned}\]
\[\begin{aligned}
r_{98}&=-x_{19}x_{128}+x_{131}
\end{aligned}\]
\[\begin{aligned}
r_{99}&=x_{127}-x_{129}x_{131}-x_{130}
\end{aligned}\]
\[\begin{aligned}
r_{100}&=-x_{131}x_{133}+x_{132}-x_{134}
\end{aligned}\]
\[\begin{aligned}
r_{101}&=-x_{19}x_{136}+x_{139}
\end{aligned}\]
\[\begin{aligned}
r_{102}&=x_{135}-x_{137}x_{139}-x_{138}
\end{aligned}\]
\[\begin{aligned}
r_{103}&=-x_{139}x_{141}+x_{140}-x_{142}
\end{aligned}\]
\[\begin{aligned}
r_{104}&=-x_{19}x_{109}+x_{146}
\end{aligned}\]
\[\begin{aligned}
r_{105}&=x_{143}-x_{144}x_{146}-x_{145}
\end{aligned}\]
\[\begin{aligned}
r_{106}&=-x_{117}x_{148}+x_{147}-x_{149}
\end{aligned}\]
\[\begin{aligned}
r_{107}&=7+x_{150}-2x_{150}x_{152}-14x_{152}+x_{153}
\end{aligned}\]
\[\begin{aligned}
r_{108}&=-x_{150}x_{153}-8x_{153}+x_{154}
\end{aligned}\]
\[\begin{aligned}
r_{109}&=2x_{151}+x_{151}^{2}+x_{153}^{2}-x_{154}^{2}
\end{aligned}\]
\[\begin{aligned}
r_{110}&=-x_{151}x_{156}-x_{156}+x_{157}
\end{aligned}\]
\[\begin{aligned}
r_{111}&=2x_{155}+x_{155}^{2}+x_{156}^{2}-x_{157}^{2}
\end{aligned}\]
\[\begin{aligned}
r_{112}&=x_{4}^{2}-2x_{4}x_{151}-2x_{151}+x_{158}
\end{aligned}\]
\[\begin{aligned}
r_{113}&=-1-x_{4}x_{156}+x_{6}-x_{155}-x_{156}+x_{157}+x_{158}x_{159}
\end{aligned}\]
\[\begin{aligned}
r_{114}&=1-2x_{5}+x_{5}^{2}-2x_{5}x_{151}+x_{160}
\end{aligned}\]
\[\begin{aligned}
r_{115}&=-x_{5}x_{206}+x_{19}+x_{160}x_{161}-x_{205}+x_{207}
\end{aligned}\]
\[\begin{aligned}
r_{116}&=-1+2x_{151}x_{221}-x_{162}+2x_{221}-x_{221}^{2}
\end{aligned}\]
\[\begin{aligned}
r_{117}&=-1+x_{107}+x_{162}x_{163}-x_{205}-x_{206}x_{221}+x_{207}
\end{aligned}\]
\[\begin{aligned}
r_{118}&=-x_{151}x_{165}-x_{165}+x_{166}
\end{aligned}\]
\[\begin{aligned}
r_{119}&=2x_{164}+x_{164}^{2}+x_{165}^{2}-x_{166}^{2}
\end{aligned}\]
\[\begin{aligned}
r_{120}&=1-2x_{107}+x_{107}^{2}-2x_{107}x_{151}+x_{167}
\end{aligned}\]
\[\begin{aligned}
r_{121}&=-1-x_{107}x_{165}+x_{108}-x_{164}+x_{166}+x_{167}x_{168}
\end{aligned}\]
\[\begin{aligned}
r_{122}&=4-4x_{107}+x_{107}^{2}-2x_{107}x_{151}+2x_{151}+x_{169}
\end{aligned}\]
\[\begin{aligned}
r_{123}&=-1-x_{107}x_{156}+x_{109}-x_{155}+x_{156}+x_{157}+x_{169}x_{170}
\end{aligned}\]
\[\begin{aligned}
r_{124}&=1+x_{151}-x_{151}x_{172}-x_{172}+x_{173}
\end{aligned}\]
\[\begin{aligned}
r_{125}&=1+2x_{171}+x_{171}^{2}-2x_{172}+x_{172}^{2}-x_{173}^{2}
\end{aligned}\]
\[\begin{aligned}
r_{126}&=-1+x_{107}-x_{107}x_{172}+x_{113}+x_{167}x_{174}-x_{171}+x_{173}
\end{aligned}\]
\[\begin{aligned}
r_{127}&=-2x_{151}-x_{151}^{2}+x_{175}
\end{aligned}\]
\[\begin{aligned}
r_{128}&=1-x_{107}-x_{107}x_{155}-x_{107}x_{157}+x_{114}+x_{151}x_{157}+x_{155}-x_{156}x_{175}+2x_{157}\\
&{}+x_{169}x_{176}
\end{aligned}\]
\[\begin{aligned}
r_{129}&=-x_{151}x_{178}-x_{178}+x_{179}
\end{aligned}\]
\[\begin{aligned}
r_{130}&=2x_{177}+x_{177}^{2}+x_{178}^{2}-x_{179}^{2}
\end{aligned}\]
\[\begin{aligned}
r_{131}&=-1-x_{107}x_{178}+x_{118}+x_{167}x_{180}-x_{177}+x_{179}
\end{aligned}\]
\[\begin{aligned}
r_{132}&=-x_{151}x_{182}-x_{182}+x_{183}
\end{aligned}\]
\[\begin{aligned}
r_{133}&=2x_{181}+x_{181}^{2}+x_{182}^{2}-x_{183}^{2}
\end{aligned}\]
\[\begin{aligned}
r_{134}&=-1-x_{107}x_{182}+x_{121}+x_{167}x_{184}-x_{181}+x_{183}
\end{aligned}\]
\[\begin{aligned}
r_{135}&=-x_{151}x_{186}-x_{186}+x_{187}
\end{aligned}\]
\[\begin{aligned}
r_{136}&=2x_{185}+x_{185}^{2}+x_{186}^{2}-x_{187}^{2}
\end{aligned}\]
\[\begin{aligned}
r_{137}&=-1-x_{107}x_{186}+x_{124}+x_{167}x_{188}-x_{185}+x_{187}
\end{aligned}\]
\[\begin{aligned}
r_{138}&=-x_{151}x_{190}-x_{190}+x_{191}
\end{aligned}\]
\[\begin{aligned}
r_{139}&=2x_{189}+x_{189}^{2}+x_{190}^{2}-x_{191}^{2}
\end{aligned}\]
\[\begin{aligned}
r_{140}&=-1-x_{107}x_{190}+x_{127}+x_{167}x_{192}-x_{189}+x_{191}
\end{aligned}\]
\[\begin{aligned}
r_{141}&=-x_{151}x_{194}-x_{194}+x_{195}
\end{aligned}\]
\[\begin{aligned}
r_{142}&=2x_{193}+x_{193}^{2}+x_{194}^{2}-x_{195}^{2}
\end{aligned}\]
\[\begin{aligned}
r_{143}&=-1-x_{107}x_{194}+x_{128}+x_{169}x_{196}-x_{193}+x_{194}+x_{195}
\end{aligned}\]
\[\begin{aligned}
r_{144}&=-x_{151}x_{198}-x_{198}+x_{199}
\end{aligned}\]
\[\begin{aligned}
r_{145}&=2x_{197}+x_{197}^{2}+x_{198}^{2}-x_{199}^{2}
\end{aligned}\]
\[\begin{aligned}
r_{146}&=-1-x_{107}x_{198}+x_{132}+x_{167}x_{200}-x_{197}+x_{199}
\end{aligned}\]
\[\begin{aligned}
r_{147}&=-x_{151}x_{202}-x_{202}+x_{203}
\end{aligned}\]
\[\begin{aligned}
r_{148}&=2x_{201}+x_{201}^{2}+x_{202}^{2}-x_{203}^{2}
\end{aligned}\]
\[\begin{aligned}
r_{149}&=-1-x_{107}x_{202}+x_{135}+x_{167}x_{204}-x_{201}+x_{203}
\end{aligned}\]
\[\begin{aligned}
r_{150}&=1-x_{151}x_{206}-x_{206}+x_{207}
\end{aligned}\]
\[\begin{aligned}
r_{151}&=-1+2x_{205}+x_{205}^{2}+x_{206}^{2}-2x_{207}-x_{207}^{2}
\end{aligned}\]
\[\begin{aligned}
r_{152}&=-x_{107}x_{206}+x_{136}+x_{169}x_{208}-x_{205}+x_{206}+x_{207}
\end{aligned}\]
\[\begin{aligned}
r_{153}&=-x_{151}x_{210}-x_{210}+x_{211}
\end{aligned}\]
\[\begin{aligned}
r_{154}&=2x_{209}+x_{209}^{2}+x_{210}^{2}-x_{211}^{2}
\end{aligned}\]
\[\begin{aligned}
r_{155}&=-1-x_{107}x_{210}+x_{140}+x_{167}x_{212}-x_{209}+x_{211}
\end{aligned}\]
\[\begin{aligned}
r_{156}&=-x_{151}x_{214}-x_{214}+x_{215}
\end{aligned}\]
\[\begin{aligned}
r_{157}&=2x_{213}+x_{213}^{2}+x_{214}^{2}-x_{215}^{2}
\end{aligned}\]
\[\begin{aligned}
r_{158}&=-1-x_{107}x_{214}+x_{143}+x_{167}x_{216}-x_{213}+x_{215}
\end{aligned}\]
\[\begin{aligned}
r_{159}&=-x_{151}x_{218}-x_{218}+x_{219}
\end{aligned}\]
\[\begin{aligned}
r_{160}&=2x_{217}+x_{217}^{2}+x_{218}^{2}-x_{219}^{2}
\end{aligned}\]
\[\begin{aligned}
r_{161}&=-1-x_{107}x_{218}+x_{147}+x_{167}x_{220}-x_{217}+x_{219}
\end{aligned}\]
\[\begin{aligned}
r_{162}&=1-2x_{0}-3x_{0}^{2}-2x_{0}x_{4}-x_{1}^{2}-x_{2}^{2}-2x_{4}-2x_{4}^{2}-x_{5}^{2}\\
&{}-x_{6}^{2}+2x_{10}-x_{10}^{2}-x_{11}^{2}-x_{12}^{2}-x_{13}^{2}-x_{14}^{2}-x_{15}^{2}\\
&{}-x_{16}^{2}-x_{17}^{2}-x_{18}^{2}-x_{19}^{2}-x_{105}^{2}+2x_{107}-2x_{107}^{2}-x_{108}^{2}\\
&{}-x_{109}^{2}-x_{113}^{2}-x_{114}^{2}-x_{118}^{2}-x_{121}^{2}-x_{124}^{2}-x_{127}^{2}\\
&{}-x_{128}^{2}-x_{132}^{2}-x_{135}^{2}-x_{136}^{2}-x_{140}^{2}-x_{143}^{2}-x_{147}^{2}+x_{150}\\
&{}-x_{221}^{2}
\end{aligned}\]
\[\begin{aligned}
r_{163}&=-x_{5}x_{10}-x_{5}x_{11}-x_{5}x_{12}-x_{5}x_{13}-x_{5}x_{14}-x_{5}x_{15}-x_{5}x_{16}\\
&{}-x_{5}x_{17}-x_{5}x_{18}+x_{221}
\end{aligned}\]
\[\begin{aligned}
r_{164}&=-16+x_{0}-x_{222}^{2}
\end{aligned}\]
\[\begin{aligned}
r_{165}&=-1+x_{2}+x_{6}-x_{6}^{2}
\end{aligned}\]
\[\begin{aligned}
r_{166}&=-x_{151}x_{224}-x_{224}+x_{233}
\end{aligned}\]
\[\begin{aligned}
r_{167}&=-1+x_{223}^{2}+x_{224}^{2}-x_{233}^{2}
\end{aligned}\]
\[\begin{aligned}
r_{168}&=-x_{226}-x_{226}x_{227}+x_{234}
\end{aligned}\]
\[\begin{aligned}
r_{169}&=-1+x_{225}^{2}+x_{226}^{2}-x_{234}^{2}
\end{aligned}\]
\[\begin{aligned}
r_{170}&=-x_{218}^{2}+x_{235}
\end{aligned}\]
\[\begin{aligned}
r_{171}&=x_{224}-2x_{229}x_{235}-x_{235}
\end{aligned}\]
\[\begin{aligned}
r_{172}&=-4x_{224}-8x_{224}x_{228}+x_{227}
\end{aligned}\]
\[\begin{aligned}
r_{173}&=-x_{151}-x_{223}x_{230}+x_{227}
\end{aligned}\]
\[\begin{aligned}
r_{174}&=-1-x_{217}-x_{223}x_{231}+x_{225}
\end{aligned}\]
\[\begin{aligned}
r_{175}&=-x_{105}-4x_{218}x_{232}+x_{226}
\end{aligned}\]
\[\begin{aligned}
r_{176}&=2x_{0}^{2}+2x_{0}x_{4}-2x_{0}x_{156}-2x_{0}x_{165}+x_{1}^{2}-2x_{1}x_{206}+x_{2}^{2}\\
&{}-2x_{2}x_{194}+x_{4}^{2}-2x_{4}x_{165}+x_{10}^{2}-2x_{10}x_{172}+x_{11}^{2}-2x_{11}x_{178}\\
&{}+x_{12}^{2}-2x_{12}x_{182}+x_{13}^{2}-2x_{13}x_{186}+x_{14}^{2}-2x_{14}x_{190}+x_{15}^{2}\\
&{}-2x_{15}x_{198}+x_{16}^{2}-2x_{16}x_{202}+x_{17}^{2}-2x_{17}x_{210}+x_{18}^{2}-2x_{18}x_{214}\\
&{}+x_{105}^{2}-2x_{105}x_{218}-x_{151}x_{236}+x_{155}^{2}+x_{156}^{2}+x_{164}^{2}+x_{165}^{2}\\
&{}+x_{171}^{2}+x_{172}^{2}+x_{177}^{2}+x_{178}^{2}+x_{181}^{2}+x_{182}^{2}+x_{185}^{2}\\
&{}+x_{186}^{2}+x_{189}^{2}+x_{190}^{2}+x_{193}^{2}+x_{194}^{2}+x_{197}^{2}+x_{198}^{2}\\
&{}+x_{201}^{2}+x_{202}^{2}+x_{205}^{2}+x_{206}^{2}+x_{209}^{2}+x_{210}^{2}+x_{213}^{2}\\
&{}+x_{214}^{2}+x_{217}^{2}+x_{218}^{2}
\end{aligned}\]
\[\begin{aligned}
r_{177}&=-1-x_{3}x_{5}+x_{5}+x_{6}-x_{237}^{2}-x_{238}^{2}-x_{239}-x_{239}^{2}
\end{aligned}\]
\[\begin{aligned}
r_{178}&=x_{3}x_{5}-x_{6}-x_{240}^{2}-x_{241}^{2}-x_{242}-x_{242}^{2}
\end{aligned}\]
\[\begin{aligned}
r_{179}&=-2-x_{105}^{2}-x_{156}^{2}-x_{165}^{2}+2x_{172}-x_{172}^{2}-x_{178}^{2}-x_{182}^{2}\\
&{}-x_{186}^{2}-x_{190}^{2}-x_{194}^{2}-x_{198}^{2}-x_{202}^{2}-x_{206}^{2}-x_{210}^{2}\\
&{}-x_{214}^{2}+x_{218}^{2}-x_{243}^{2}-x_{244}^{2}-x_{245}-x_{245}^{2}
\end{aligned}\]
\[\begin{aligned}
r_{180}&=-2+x_{107}-4x_{246}-4x_{246}^{2}
\end{aligned}\]
\[\begin{aligned}
r_{181}&=-x_{107}x_{111}+x_{111}+x_{247}
\end{aligned}\]
\[\begin{aligned}
r_{182}&=-x_{107}x_{145}+x_{145}+x_{248}
\end{aligned}\]
\[\begin{aligned}
r_{183}&=-1+x_{114}^{2}-4x_{116}^{2}-4x_{120}^{2}-4x_{123}^{2}-4x_{126}^{2}-4x_{149}^{2}-4x_{247}^{2}\\
&{}-4x_{248}^{2}-x_{249}^{2}-x_{250}^{2}-x_{251}-x_{251}^{2}
\end{aligned}\]
\[\begin{aligned}
r_{184}&=-1+x_{128}^{2}-4x_{130}^{2}-4x_{134}^{2}-x_{252}^{2}-x_{253}^{2}-x_{254}-x_{254}^{2}
\end{aligned}\]
\[\begin{aligned}
r_{185}&=-1+x_{136}^{2}-4x_{138}^{2}-4x_{142}^{2}-x_{255}^{2}-x_{256}^{2}-x_{257}-x_{257}^{2}
\end{aligned}\]

\endgroup
\end{document}
